\section{Question 2}
% (1) L1 の要素数
\begin{frame}
  \frametitle{Question 2: (1) $L_1$ の要素数}
  
  \textbf{問題:} $L_1$ の要素数を答えなさい。
  
  $L_1$ は $\Sigma = \{a, b, c\}$ 上の長さ 5 以下のすべての文字列の集合である。
  
  \begin{alertblock}{解答}
    空文字列も含めて、
    $|L_1| = 3^0 + 3^1 + 3^2 + 3^3 + 3^4 + 3^5 = 1 + 3 + 9 + 27 + 81 + 243 = 364$
    である。
  \end{alertblock}

  \begin{itemize}
    \item $\Sigma = \{a, b, c\}$ なので、各文字は3通り。
    \item 長さ0（空文字列）から長さ5まで、各長さごとに $3^n$ 通り（$n$は長さ）。
    \item それぞれを合計して364個となる。
  \end{itemize}
\end{frame}

% (2) L2 の要素数

% (2-1) L2 の要素数（考え方と概要）
\begin{frame}
  \frametitle{Question 2: (2) $L_2$ の要素数（1/2）}
  \textbf{問題:} $L_2$ の要素数を答えなさい。
  
  $L_2$ は $L_1$ のうち、部分文字列として ``aba'' を含むものの集合である。
  
  \begin{alertblock}{考え方}
    $L_1$ の全要素から ``aba'' を含まないものを除外することで求める。
    すなわち、$|L_2| = |L_1| - $（``aba'' を含まない文字列の個数）
  \end{alertblock}

  \begin{itemize}
    \item $L_1$ の要素数は $364$ 個（長さ0～5の全ての文字列）。
    \item $L_2$ は ``aba'' をどこかに含むもの。
    \item 例えば ``aba'', ``ababa'', ``caba'', ``abac'' など。
    \item ``aba'' を含まないものを全体から引くことで個数を求める。
  \end{itemize}
\end{frame}

% (2-2) L2 の要素数（詳細な計算・具体例）
\begin{frame}
  \frametitle{Question 2: (2) $L_2$ の要素数（2/2）}
  \textbf{詳細な計算（具体的な数え上げ）}

  \begin{itemize}
    \item \textbf{長さ3:} ``aba'' だけなので $1$ 個
    \item \textbf{長さ4:}  
      \begin{itemize}
        \item 先頭に ``aba''：\ \texttt{aba*} の形 $\to$ 3通り
        \item 後ろに ``aba''：\ \texttt{*aba} の形 $\to$ 3通り
        \item 合計 $6$ 個
      \end{itemize}
    \item \textbf{長さ5:}
      \begin{itemize}
        \item 先頭に ``aba''：\ \texttt{aba**} の形 $3 \times 3 = 9$
        \item 真ん中に ``aba''：\ \texttt{*aba*} の形 $3 \times 3 = 9$
        \item 後ろに ``aba''：\ \texttt{**aba} の形 $3 \times 3 = 9$
        \item ただし ``ababa'' は重複するので $-1$
        \item 合計 $9 + 9 + 9 - 1 = 26$
      \end{itemize}
    \begin{alertblock}{解答}
      \begin{center}
        $1 + 6 + 26 = 33$ 個
      \end{center}
    \end{alertblock}
  \end{itemize}
  \vspace{1em}
\end{frame}

% (3) L3 の辞書式順序で最初の10要素
\begin{frame}
  \frametitle{Question 2: (3) $L_3$ の辞書式順序で最初の10要素}
  \textbf{問題:} $L_3$ の要素のうち辞書式順序で最初の10要素を列挙しなさい。
  
  $L_3$ は $L_1$ のうち回文であるものの集合である。

  \begin{alertblock}{解答}
      \begin{center}
        \varepsilon, a, b, c, aa, bb, cc, aaa, aba, aca
      \end{center}
  \end{alertblock}

  \begin{itemize}
    \item 回文とは、前から読んでも後ろから読んでも同じ文字列。
    \item 例えば ``aba'' や ``bab''、``ccbcc'' など。
    \item 辞書式順序で並べている。
    \item 空文字列（$\varepsilon$）も回文として含まれる。
    \item 短いものから順に並ぶ。
  \end{itemize}
\end{frame}

% (4) L3 の辞書式順序で最後の5要素
\begin{frame}
  \frametitle{Question 2: (4) $L_3$ の辞書式順序で最後の5要素}
  \textbf{問題:} $L_3$ の要素のうち辞書式順序で最後の5要素を列挙しなさい。
  
  \begin{alertblock}{解答}
      \begin{center}
        cbbbc,\ cbcbc,\ ccacc,\ ccbcc,\ ccccc
      \end{center}
  \end{alertblock}

  \begin{itemize}
    \item アルファベット順で最後に来る回文を5つ並べている。
    \item 長さが大きいほど後ろに並ぶ。
  \end{itemize}
\end{frame}